\item Está trabajando como testigo experto para el propietario de un complejo de rascacielos en el centro de la ciudad. El propietario está siendo demandado por peatones en las calles debajo de sus edificios que resultaron heridos por la caída de cristales cuando ventanas salieron expelidas de los lados del edificio. El efecto Bernoulli puede tener consecuencias importantes para las ventanas en dichos edificios. Por ejemplo, el viento puede soplar alrededor de un rascacielos a una velocidad notablemente alta, creando una baja presión en la superficie exterior de las ventanas. La presión atmosférica más alta en el aire inmóvil dentro de los edificios puede hacer que las ventanas sean expulsadas. (a) En su investigación sobre el caso, encontrará algunas vistas generales del proyecto de su cliente, como se muestra en la figura P14.26. El proyecto incluye dos altos rascacielos y un área de parque en un terreno cuadrado. El plan (i) fue presentado por los arquitectos y planificadores originales. En el último minuto, el propietario decidió que no quería que los terrenos del parque se dividieran en dos áreas y presentó el Plan (ii), que es la forma en que se construyó el proyecto. Explique a su cliente por qué el plan (ii) es una situación mucho más peligrosa en términos de ventanas expulsadas que el plan (i). (b) Su cliente no está convencido por su argumento conceptual en el inciso (a), por lo que proporciona un argumento numérico. Un viento horizontal sopla con una velocidad de $11.2 \text{ m/s}$ fuera de un gran panel de vidrio con dimensiones de $4.00 \text{ m} \times 1.50 \text{ m}$. Suponga que la densidad del aire es constante a $1.20 \text{ kg/m}^3$. El aire dentro del edificio está a presión atmosférica. Calcule la fuerza total ejercida por el aire sobre el cristal de la ventana para su cliente. (c) ¿Qué pasaría si? Para convencer aún más a su cliente de los problemas con el diseño del edificio, calcule la fuerza total ejercida por el aire sobre el cristal si la velocidad del viento en los edificios es de $22.4 \text{ m/s}$, dos veces más alta que en el inciso (b).
